% Chapter 5 -- Using Applications (Part I, Learning; AN-14 task form).
% Sources (Law 1): os/milton/kmain.c finder_surface_open (an icon whose name
% ends .EXE is launched; any other file is not opened from the Finder),
% flair_live_launch_app / tbx_launch, flair_live_sync_foreground (the
% foreground application's bar in band 2; a click on the desktop makes the
% Finder foreground), sf_run (Standard File: the folder pop-up, Desktop,
% Eject and New dimmed, the name field, sf_name_ok: 8.3 names, typed letters
% shown in capitals, the Replace question); os/flair/tbxgate.c (one disk
% application at a time: TENANT-SLOT-BUSY); os/apps/itext/itext.c
% (InitechText: the untitled window, Save on an untitled document asks for a
% name, Quit with unsaved changes asks Save / Cancel / Don't Save);
% os/apps/calcda/calcda.c and os/flair/apple_menu.c (the Calculator from
% the Apple menu); kmain.c finder_power / finder_exit_dos (an open
% application is asked to quit).  Model: Macintosh User's Guide 1991 Ch. 3-5
% "Setting Up Your Programs", "Creating a Document and Saving Your Work",
% "Changing a Document" (apple-macintosh.md Sec 8.2).
% GATES: test-flair-app-launch test-flair-itext test-itext-route test-stdfile test-flair-calcda test-flair-apple-menu test-flair-tenant-menu test-flair-appswitch test-pubs-flair-screens
\Gates{test-flair-app-launch, test-flair-itext, test-itext-route,
  test-stdfile, test-flair-calcda, test-flair-apple-menu,
  test-flair-tenant-menu, test-flair-appswitch, test-pubs-flair-screens}
\chapter{Using Applications}

An \term{application} is a program you use to do your work: to write a
memo, to keep a ledger, to work out a sum.  The information you create with
an application is a \term{document}.  This chapter shows you how to open an
application, how to work with its menus, how to save and open documents,
and how to quit.  The examples use InitechText, the text editor that comes
with FLAIR; the \doctitle{FLAIR Accessories Guide} describes it fully.

\section{Opening an application}
\index{applications!opening}\index{opening an application}
\Operation{TENANT-LOAD}
Applications are kept in the \file{APPS} folder of the data disk.  Their
icons are diamonds, and their names end in \file{.EXE}.
\begin{TaskSteps}
\Step Open the \file{APPS} folder.
\Result{To see all its icons at once, click the zoom box of its window.}
\Step Double-click \file{TEXT.EXE}.
\Result{FLAIR loads InitechText from the disk.  Its window, titled
\texttt{untitled}, opens in front of the other windows, ready for typing,
and the lower menu bar changes to InitechText's menus: the Apple menu,
\menuname{File} and \menuname{Edit}.}
\Step Type a line of text.
\Result{The text appears at the \term{insertion
point}, the thin vertical line at the top left of the window.}
\end{TaskSteps}

\FlairScreen[label=fig:textopen]{textopen}{InitechText, opened from the
APPS folder}

FLAIR runs one application from the disk at a time.  While InitechText is
open, double-clicking another application does nothing; quit InitechText
first.  HELLO and NOTES, which are built into FLAIR, are always available
alongside it.

Only applications open from the desktop.  Double-clicking a document, or a
program whose name ends in \file{.COM}, does nothing: open a document from
within its application, and run the \file{.COM} programs at the InitechDOS
prompt (Chapter~1).

\section{The active application}
\index{active application}\index{applications!switching between}
The lower menu bar always shows the menus of the active application.  To
change the active application, click one of its windows:
\begin{itemize}
\item Click InitechText's window to work in InitechText.
\item Click a Finder window, or a bare part of the desktop, to work with
  the Finder.  The Finder's windows come to the front.
\item Click HELLO or NOTES to bring that window to the front.
\end{itemize}
An application stays open when it is not active.  Its windows keep their
contents, and you return to it by clicking one of them.

\section{Saving a document}
\index{saving a document}\index{Save command}
\Operation{TENANT-SF}
Until you save it, a document exists only in the memory of the computer and
is lost when the application quits.  Saving writes it to the disk as a
file.
\begin{TaskSteps}
\Step Choose \menuitem{Save} from InitechText's \menuname{File} menu, or
press \key{Ctrl-S}.
\Result{Because the document has never been saved, a dialog box asks where
to save it and what to call it.}
\end{TaskSteps}

\FlairScreen[label=fig:textsaveas]{textsaveas}{The dialog box for saving a
document}

The dialog box has these parts:

\SideHead{Folder menu}The box at the top shows the folder the document will
be saved in; at first it is the top level of the data disk,
\texttt{INITECH}.  Click it to move up to the folder that holds this one.

\SideHead{List}The files and folders in that folder.  Double-click a
folder to save the document inside it.

\SideHead{Desktop}Click \texttt{Desktop} to go back to the top level of the
disk.  \texttt{Eject} and \texttt{New} are dimmed: the data disk cannot be
ejected while FLAIR runs, and new folders are made in the Finder.

\SideHead{Name}The box under \texttt{Save this document as:} holds the name
the document will have.  It starts as \texttt{UNTITLED}.

\begin{TaskSteps}[resume]
\Step Press \key{Backspace} until the name box is empty, then type
\type{memo.txt}.
\Result{What you type appears in capital letters, because InitechDOS
stores names in capitals.  The name must have one to eight characters,
and may be followed by a period and up to three more.}
\Step Click \texttt{Save}, or press \key{Enter}.
\Result{The document is written to the disk, and the window takes the name
of the file.}
\end{TaskSteps}

\FlairScreen[label=fig:textsaved]{textsaved}{The document saved as
MEMO.TXT}

To save the document again after more changes, choose \menuitem{Save}: it
now saves under the same name without asking.  To save a copy under
another name, choose \menuitem{Save As...}.  If you give a name that is
already used in the folder, FLAIR asks whether to replace the file that has
that name: click \texttt{Replace} to replace it, or \texttt{Cancel} to type
another name.

\section{Opening a document}
\index{opening a document}\index{Open command}
\begin{TaskSteps}
\Step Choose \menuitem{Open...} from InitechText's \menuname{File} menu,
or press \key{Ctrl-O}.
\Result{A dialog box lists the files and folders of the top level of the
data disk.}
\end{TaskSteps}

\FlairScreen[label=fig:textopenbox]{textopenbox}{The dialog box for opening
a document}

\begin{TaskSteps}[resume]
\Step Click the name of the document you want, then click \texttt{Open}.
\Result{You can also double-click the name, or move the highlight with
\key{Up Arrow} and \key{Down Arrow} and press \key{Enter}.  To look inside a
folder, double-click it.  The document opens in InitechText's window.}
\end{TaskSteps}

\section{Quitting an application}
\index{quitting}\index{Quit command}
When you have finished with an application, quit it to free the computer
for another.
\begin{TaskSteps}
\Step Choose \menuitem{Quit} from the application's \menuname{File} menu,
or press \key{Ctrl-Q}.
\Result{If the document has changes you have not saved, the application
asks what to do with them.}
\end{TaskSteps}

\FlairScreen[label=fig:textask]{textask}{InitechText asks about unsaved
changes}
\Operation{TENANT-ALERT}

\SideHead{Save}Saves the document (asking for a name if it has none) and
then quits.  \texttt{Save} is the default button: pressing \key{Enter}
chooses it.
\SideHead{Don't Save}Quits without saving; the changes are lost.
\SideHead{Cancel}Goes back to the document; nothing is quit.  Pressing
\key{Esc} chooses it.

When you choose \menuitem{Restart}, \menuitem{Shut Down} or \menuitem{Exit
to InitechDOS}, FLAIR asks the open application to quit in the same way.
If you click \texttt{Cancel} in its alert, the computer is not restarted or
shut down.

\section{Opening the Calculator}
\index{Calculator}
Some small applications, called \term{accessories}, are listed in the
Apple menu so that you can open them from any application.
\Operation{APPLE-LAUNCH}
\begin{TaskSteps}
\Step Choose \menuitem{Calculator} from the Apple menu.
\Result{The Calculator opens.  It works like a pocket calculator: click
its keys, or type the numbers and signs on the keyboard.  The
\doctitle{FLAIR Accessories Guide} describes it.}
\end{TaskSteps}

\FlairScreen[label=fig:calcda]{calcda}{The Calculator}

The Calculator is a disk application like the others, so it does not open
while another application from the disk is open.  To close it, click its
close box, or choose \menuitem{Close} from its \menuname{File} menu.
